\section{Finite rows and their scalar products}\label{sec:finite}

The first step is finite-dimensional.  We express every fixed word of
marked rows as a finite sum of determinants on the Bethe roots, with all
normalizing factors specified.  The construction of those roots and the
estimates needed to pass to infinite volume are deferred to
Section~\ref{sec:thermodynamic}.  In particular, no completeness assertion
for Bethe eigenvectors is used in this section.

\subsection{The physical row and the standard monodromy}

The algebra below holds for $0<\lambda<\pi$, corresponding to the
open parameter interval $0<c<2$. The endpoint $c=2$ is treated
separately in Appendix~\ref{sec:endpoint-rotation}. Set
\begin{equation}\label{eq:spectral-coordinates}
 \lambda=\arccos(-\Delta),\qquad \eta=\lambda/\pi,\qquad
 v=\pi/2,\qquad x=i\eta z,\qquad w=e^{iz}.
\end{equation}
The physical time angle is $z=0$, increasing time points east, and the
space angle is $v$, increasing space points north.  Thus a common
imaginary translation of the $z$ arguments is a common real translation
of their $x$ arguments.  At a path step, the spin is $+1$ when the arrow
points along its left normal; the height increment is minus that spin.
These conventions fix the sign of every marked row.

Let $L$ be even and let $b_0,\ldots,b_{L-1}$ be real numbers.  On
$\mathcal H_L=(\C^2)^{\otimes L}$ use the orthonormal spin basis, and put
\[
 Z=\begin{pmatrix}1&0\\0&-1\end{pmatrix},\qquad
 X=\begin{pmatrix}0&1\\1&0\end{pmatrix}.
\]
An auxiliary copy of $\C^2$ is denoted by the subscript $a$.  Products
over the sites below have increasing site index from left to right.  Define
\begin{align}
 R(u)&=
 \begin{pmatrix}
 \sin(\lambda-u)/\sin\lambda&0&0&0\\
 0&\sin u/\sin\lambda&1&0\\
 0&1&\sin u/\sin\lambda&0\\
 0&0&0&\sin(\lambda-u)/\sin\lambda
 \end{pmatrix},\label{eq:physical-R}\\
 M_{\rm ph}(x)&=\prod_{j=0}^{L-1}
 R_{aj}\bigl(\lambda/2+i(x-b_j)\bigr),\qquad
 T^D(x)=\tr_a\bigl(D_aM_{\rm ph}(x)\bigr),\label{eq:physical-row}
\end{align}
where $D$ is an arbitrary complex $2\times2$ matrix.  Its row index is
on the right tile of the horizontal seam and its column index on the
left.  The four indices in \eqref{eq:physical-R} send the top and left
step spins to the bottom and right step spins.  When $b_j=0$ and $x=0$,
dividing all vertex weights by $\sin(\lambda/2)/\sin\lambda$ gives
$a=b=1$ and $c=2\cos(\lambda/2)$, hence $\Delta=-\cos\lambda$.

Write $[u]=\sinh u$ and $s_\lambda=[i\lambda]=i\sin\lambda$.
The standard monodromy used in the algebraic Bethe ansatz is
\begin{align}
 L_{aj}(x-b_j)&=
 \begin{pmatrix}
 [x-b_j+i\lambda/2]/s_\lambda&0&0&0\\
 0&[x-b_j-i\lambda/2]/s_\lambda&1&0\\
 0&1&[x-b_j-i\lambda/2]/s_\lambda&0\\
 0&0&0&[x-b_j+i\lambda/2]/s_\lambda
 \end{pmatrix},\notag\\
 M(x)&=\prod_{j=0}^{L-1}L_{aj}(x-b_j)
      =\begin{pmatrix}\mathsf A(x)&\mathsf B(x)\\
                         \mathsf C(x)&\mathsf L(x)\end{pmatrix}_a.
 \label{eq:standard-monodromy}
\end{align}
In particular, the physical mixed diagonal entry has the opposite sign.
Let
\begin{equation}\label{eq:parity-gauges}
 \mathcal P_L=\prod_{j\ {\mathrm{even}}}Z_j,
 \qquad \mathcal Q_L=\prod_{j=0}^{L-1}Z_j.
\end{equation}

\begin{lemma}[Finite row algebra]\label{lem:row-algebra}
For even $L$, real $b_j$, and all complex spectral arguments,
\begin{equation}\label{eq:monodromy-gauge}
 M_{\rm ph}(x)=\mathcal P_L M(x)\mathcal P_L\mathcal Q_L,
 \qquad
 T^D(x)^*=T^{X\overline D X}(\bar x).
\end{equation}
The rows with a common diagonal twist commute.  More generally, put
\[
 \mathcal R(u)=
 \begin{pmatrix}
 [u+i\lambda]/s_\lambda&0&0&0\\
 0&[u]/s_\lambda&1&0\\
 0&1&[u]/s_\lambda&0\\
 0&0&0&[u+i\lambda]/s_\lambda
 \end{pmatrix},\qquad
 \mathcal R_{\rm ph}(u)=Z_a\mathcal R(u)Z_b.
\]
Then
\begin{equation}\label{eq:physical-RTT}
 \mathcal R_{{\rm ph},ab}(p-q)M_{{\rm ph},a}(p)M_{{\rm ph},b}(q)
 =M_{{\rm ph},b}(q)M_{{\rm ph},a}(p)
   \mathcal R_{{\rm ph},ab}(p-q).
\end{equation}
Where the intertwiner is invertible, exchanging two rows conjugates
their joint auxiliary mark by $\mathcal R_{\rm ph}(p-q)$.
For the homogeneous chain, $T^I(i\lambda/2)$ is the cyclic spin
translation and $T^I(-i\lambda/2)$ is its inverse.
\end{lemma}

\begin{proof}
The local identity is
$R_{aj}(\lambda/2+i(x-b_j))=Z_aL_{aj}(x-b_j)Z_j$.
Both local matrices commute with $Z_aZ_j$.  Moving the auxiliary $Z$'s
through the product, and using that $L$ is even, gives
\eqref{eq:monodromy-gauge}.  This also proves that the gauge is
independent of the roots, spectral arguments, and inhomogeneities.

If $N_{\alpha\beta}(x)$ is one auxiliary block of a physical local
matrix, direct conjugation of its entries gives
$N_{\alpha\beta}(x)^*=N_{\bar\alpha\bar\beta}(\bar x)$, where a bar on
a spin index exchanges its two values.  In an entry of a monodromy,
different local factors act on different quantum sites.  They therefore
commute when the adjoint reverses their order.  Consequently
\[
 (M_{{\rm ph},\alpha\beta}(x))^*
 =M_{{\rm ph},\bar\alpha\bar\beta}(\bar x).
\]
Taking the marked trace proves the second identity in
\eqref{eq:monodromy-gauge}; no reversal of the spatial inhomogeneities
is involved.

The standard six-vertex Yang--Baxter identity gives
$\mathcal R_{ab}(p-q)M_a(p)M_b(q)
=M_b(q)M_a(p)\mathcal R_{ab}(p-q)$.
It is a local identity of the three displayed six-vertex matrices:
spin conservation reduces its nonconstant entries to the addition
formula for $\sinh$.  Moreover,
\[
 M_{{\rm ph},a}(p)M_{{\rm ph},b}(q)
 =Z_b\mathcal P_L M_a(p)M_b(q)\mathcal P_L Z_b,
\]
and the analogous reversed product has $Z_a$ in place of $Z_b$.
Substitution yields \eqref{eq:physical-RTT}.  Taking auxiliary traces
proves the assertion about marks.  For a common diagonal matrix $D$,
$D\otimes D$ commutes with the intertwiner, so the two twisted rows
commute.  Finally $R(0)$ is the permutation matrix, and tracing a chain
of such permutations gives the cyclic translation.  The other endpoint
follows by the adjoint identity.
\end{proof}

\subsection{Bethe vectors, counting functions, and finite regularity}

Let $|0\rangle=|+\rangle^{\otimes L}$.  It is annihilated by
$\mathsf C(x)$, and its two diagonal eigenvalues are
\begin{equation}\label{eq:vacuum-functions}
 a(x)=\prod_{j=0}^{L-1}\frac{[x-b_j+i\lambda/2]}{s_\lambda},
 \qquad
 l(x)=\prod_{j=0}^{L-1}\frac{[x-b_j-i\lambda/2]}{s_\lambda}.
\end{equation}
The scalar functions in the creation-entry commutation rules are
\begin{equation}\label{eq:fgbeta}
 f(u)=\frac{[u+i\lambda]}{[u]},\qquad
 g(u)=\frac{s_\lambda}{[u]},\qquad
 \beta(u)=\frac{g(u)}{f(u)}=\frac{s_\lambda}{[u+i\lambda]}.
\end{equation}
In particular $\beta(0)=1$ and $1/f(0)=0$ by continuation.
For a finite list $U$ set
$|U\rangle=\prod_{u\in U}\mathsf B(u)|0\rangle$.
The creation entries commute, so this vector is symmetric in $U$.
Fix an integer $0\le n\le L$.  For a list
$Y=(y_1,\ldots,y_n)$ of distinct roots define
\begin{align}
 F_Y(x)&=\prod_{y\in Y}\frac{[y-x+i\lambda]}{[y-x]},\notag\\
 q_Y(x)&=\frac{l(x)}{a(x)}
          \prod_{y\in Y}\frac{[x-y+i\lambda]}{[x-y-i\lambda]},
 \label{eq:bethe-functions}\\
 \tau_t(x;Y)&=e^{it}a(x)F_Y(x)
                   \bigl(1+e^{-2it}q_Y(x)\bigr).
 \label{eq:bethe-eigenvalue}
\end{align}
Here $t$ is real near zero.  We say that $Y$ is on shell for $t$ when
\begin{equation}\label{eq:bethe-equations}
 q_Y(y_j)=-e^{2it}\qquad(1\le j\le n).
\end{equation}
The apparently singular eigenvalue in \eqref{eq:bethe-eigenvalue} then
has removable poles at the roots.

For real $u$ put
\begin{equation}\label{eq:bare-kernels}
 s(u)=\frac{\sin\lambda}{\pi(\cosh(2u)-\cos\lambda)},\qquad
 K(u)=\frac{\sin(2\lambda)}
               {\pi(\cosh(2u)-\cos(2\lambda))}.
\end{equation}
At $\lambda=\pi/2$, the second function is identically zero.  Let $S$
and $\Theta$ be their odd primitives, vanishing at zero.  Define
\begin{equation}\label{eq:counting-function}
 \xi_Y(u)=\sum_{j=0}^{L-1}S(u-b_j)-\sum_{k=1}^n\Theta(u-y_k).
\end{equation}
Consecutive real roots are solutions of
\begin{equation}\label{eq:consecutive-counting}
 \xi_Y(y_j)=j-\frac{n+1}{2}+\frac t\pi,
 \qquad y_1<\cdots<y_n.
\end{equation}

\begin{lemma}[Finite Bethe data]\label{lem:bethe-data}
Let $L$ be even, $b_j$ real, and $0<\lambda<\pi$.
A real solution of \eqref{eq:consecutive-counting} satisfies
\eqref{eq:bethe-equations}.  If its Bethe vector is nonzero, then
$|Y\rangle$ is an eigenvector of
$e^{it}\mathsf A(x)+e^{-it}\mathsf L(x)$ with eigenvalue
$\tau_t(x;Y)$, and $\mathcal P_L|Y\rangle$ is an eigenvector of the
physical twisted row with eigenvalue $(-1)^n\tau_t(x;Y)$.
The Jacobian of \eqref{eq:consecutive-counting} is
\begin{equation}\label{eq:Gaudin-matrix}
 G_{jk}=\xi_Y'(y_j)\delta_{jk}+K(y_j-y_k),\qquad
 \xi_Y'(u)=\sum_{a=0}^{L-1}s(u-b_a)-\sum_{k=1}^nK(u-y_k).
\end{equation}
If $\det G\ne0$, the roots have a real analytic continuation for real
$t$ near the given value, and on that branch
\begin{equation}\label{eq:twist-root-derivative}
 G\frac{\dd Y}{\dd t}=\frac1\pi\bbone.
\end{equation}
For $\lambda\ge\pi/2$, $G$ is positive definite at every real root
list.  Whenever $\xi_Y'(y_j)>0$, write
\begin{equation}\label{eq:Gaudin-weights}
 \ell_j=\frac1{\xi_Y'(y_j)},\qquad
 (K\ell)_{jk}=K(y_j-y_k)\ell_k;
 \qquad G\diag(\ell)=I+K\ell.
\end{equation}
These assertions concern finite matrices; they give no uniform bound on
$(I+K\ell)^{-1}$ as $L$ increases.
\end{lemma}

\begin{proof}
Logarithmic differentiation gives $q_Y'/q_Y=2\pi i\xi_Y'$.
Comparing the limits as $u\to+\infty$ fixes the integration constant:
\begin{equation}\label{eq:q-counting-phase}
 q_Y(u)=(-1)^n e^{2\pi i\xi_Y(u)}\qquad(u\in\R).
\end{equation}
Indeed $S(+\infty)=(1-\eta)/2$ and
$\Theta(+\infty)=(1-2\eta)/2$, while
$q_Y(+\infty)=e^{-i\lambda L+2i\lambda n}$ and $L$ is even.
Equation~\eqref{eq:consecutive-counting} now gives the Bethe equations.
The two diagonal action formulas proved in
Lemma~\ref{lem:finite-actions} below give the coefficient replacing a
root $p$ by $x$ as
\[
 g(x-p)\left(
 e^{it}a(p)\prod_{y\ne p}f(y-p)
 -e^{-it}l(p)\prod_{y\ne p}f(p-y)\right).
\]
The Bethe equation, with the self factor $-1$ removed, makes the
parenthesis zero.  The unchanged coefficient is
\eqref{eq:bethe-eigenvalue}.  The physical eigenvalue follows from
\eqref{eq:monodromy-gauge}, since $\mathcal Q_L=(-1)^n$ on the sector
with $n$ minus spins.

Differentiating the left side of
\eqref{eq:consecutive-counting} with respect to every root gives
\eqref{eq:Gaudin-matrix}; the $K(0)$ term cancels the derivative of the
self interaction.  The analytic implicit-function theorem gives the
local branch, and differentiation with respect to $t$ gives
\eqref{eq:twist-root-derivative}.  When $K\le0$,
\[
 G_{jj}=\sum_a s(y_j-b_a)+\sum_{k\ne j}|K(y_j-y_k)|
       >\sum_{k\ne j}|G_{jk}|.
\]
The matrix is real symmetric, so strict diagonal dominance proves
positive definiteness.  The formula for $G\diag(\ell)$ follows
entry by entry.
\end{proof}

In the remainder of the paper, a \emph{regular real on-shell list}
means a list of distinct real roots satisfying the Bethe equations,
with $\xi_Y'(y_j)>0$ for every $j$ and $\det G_Y\ne0$.
Section~\ref{sec:thermodynamic} constructs lists with these properties.

We also record the full-line kernels that will be used to compare the
finite matrices with convolution.  With
$\widehat f(p)=\int_{\R}e^{-ipu}f(u)\dd u$,
\begin{align}
 \widehat s(p)&=\frac{\sinh((\pi-\lambda)p/2)}{\sinh(\pi p/2)},&
 \widehat K(p)&=\frac{\sinh((\pi-2\lambda)p/2)}{\sinh(\pi p/2)},
 \label{eq:kernel-transforms}\\
 h:=\|K\|_{L^1(\R)}&=|1-2\eta|<1,&
 r(u):=(I+K*)^{-1}s(u)&=\frac1{2\lambda\cosh(u/\eta)}.
 \label{eq:dressed-density}
\end{align}
Here $K*$ denotes convolution with $K$.  To verify the transforms,
integrate $e^{-ipu}\sin\theta/[\pi(\cosh(2u)-\cos\theta)]$ around a
rectangle of height $i\pi$, for $0<\theta<2\pi$.  The two residues at
$i\theta/2$ and $i(\pi-\theta/2)$ give
$\sinh((\pi-\theta)p/2)/\sinh(\pi p/2)$; coincident or zero values
follow by continuity.  Take $\theta=\lambda,2\lambda$.
The kernel $K$ has constant sign, proving the $L^1$ formula, and
$\widehat s/(1+\widehat K)=1/(2\cosh(\lambda p/2))$ proves the
formula for $r$.  In particular $\int r=1/2$.
We write
\begin{equation}\label{eq:full-line-resolvent}
 R_0*=K*(I+K*)^{-1};
 \qquad R_0+K*R_0=K.
\end{equation}
The inverse exists, for example on bounded functions, by the Neumann
series since $h<1$.

\subsection{Exact actions on products of creation entries}

For the next formulas all arguments are initially distinct and avoid
the zeros of their displayed denominators.  They are identities of
meromorphic functions; coincident cases are obtained by taking the
limit of the complete expression.  We use the scalar functions
$f,g,\beta$ defined in \eqref{eq:fgbeta}.

\begin{lemma}[Single-entry actions]\label{lem:finite-actions}
For a finite spectral set $U$, put $E=U\cup\{v'\}$.  Then
\begin{align}
 \mathsf B(v')|U\rangle&=|E\rangle,\notag\\
 \mathsf A(v')|U\rangle
 &=\sum_{p\in E}a(p)\beta(v'-p)
          \prod_{u\in E\setminus\{p\}}f(u-p)
             |E\setminus\{p\}\rangle,\label{eq:A-action}\\
 \mathsf L(v')|U\rangle
 &=\sum_{p\in E}l(p)\beta(p-v')
          \prod_{u\in E\setminus\{p\}}f(p-u)
             |E\setminus\{p\}\rangle,\label{eq:L-action}\\
 \mathsf C(v')|U\rangle
 &=\sum_{\substack{p,q\in E\\p\ne q}}
       a(p)l(q)\beta(v'-p)\beta(q-v')f(q-p)
       \prod_{u\in E\setminus\{p,q\}}f(u-p)f(q-u)
             |E\setminus\{p,q\}\rangle.\label{eq:C-action}
\end{align}
The last sum is over ordered pairs: $p$ belongs to the upper diagonal
factor and $q$ to the lower diagonal factor.  Thus a word of $m$ entries
is an explicit sum over paths that add the new argument at each step
and remove zero, one, or two elements.  If the final set has the
original cardinality, its sets of removed original roots and surviving
added arguments have equal size, at most $m$.
\end{lemma}

\begin{proof}
The standard RTT relation gives commuting creation entries and
\begin{align*}
 \mathsf A(p)\mathsf B(u)
  &=f(u-p)\mathsf B(u)\mathsf A(p)
        +g(p-u)\mathsf B(p)\mathsf A(u),\\
 \mathsf L(p)\mathsf B(u)
  &=f(p-u)\mathsf B(u)\mathsf L(p)
        +g(u-p)\mathsf B(p)\mathsf L(u),\\
 [\mathsf C(p),\mathsf B(u)]
  &=g(p-u)\bigl(\mathsf A(u)\mathsf L(p)
                    -\mathsf A(p)\mathsf L(u)\bigr).
\end{align*}
Move the diagonal entry successively through the creation entries and
apply its vacuum value.  Combining terms that replace the same root
uses
\[
 f(u-h)-f(u-p)=g(p-u)\frac{[p-h]}{[u-h]}.
\]
This gives \eqref{eq:A-action}; reversing differences gives
\eqref{eq:L-action}.  These formulas include the unchanged term $p=v'$
because $\beta(0)=1$.

For completeness, the two-removal induction can be checked before any
on-shell specialization.  Induct on $|U|$. Strip one creation at $u$ from the product,
use the commutator, and apply the two diagonal formulas.  For fixed
upper and lower removals $h,k$, divide the coefficients by
\[
 a(h)l(k)f(k-h)\prod_{s\ {\mathrm{remaining}}}f(s-h)f(k-s).
\]
The required addition of the three resulting contributions is
\begin{align*}
 \beta(v'-h)\beta(k-v')
 &=\frac{\beta(v'-h)\beta(k-v')}{f(u-h)f(k-u)}\\
 &\quad+g(v'-u)\left(
       \frac{\beta(u-h)\beta(k-v')}{f(k-u)}
       -\frac{\beta(v'-h)\beta(k-u)}{f(k-v')}\right).
\end{align*}
Substituting \eqref{eq:fgbeta} and multiplying by the displayed
denominators reduces this identity to the addition formula for
$\sinh$.  The initial case is $\mathsf C(v')|0\rangle=0$; for one
creation the two ordered removals give the displayed commutator on the
vacuum.  The induction proves \eqref{eq:C-action}, including cases
where $p$ or $q$ is the newly added argument.  Finally, there are only
$m$ newly added arguments.  If the final cardinality equals the initial
one, the number of absent original roots equals the number of
surviving added arguments and is therefore at most $m$.
\end{proof}

We now have an exact expansion of a finite row word into off-shell
Bethe vectors.  The remaining finite-dimensional step is to take their
scalar products with the initial vector and cancel the vacuum factors.

\subsection{Slavnov's identity with its normalization factors}

The scalar-product determinant below specializes to the norm determinant
associated with Gaudin, McCoy and Wu and proved by
Korepin~\cite{GaudinMcCoyWu1981,Korepin1982}. We use the twisted
Slavnov identity to keep the exact prefactors in our convention visible.

Let $Y$ be on shell for $t$ and have $n$ distinct real roots.  For a
list $U$ of length $n$, define the algebraic scalar product
\begin{equation}\label{eq:ABA-scalar}
 S(Y,U)=\langle0|\prod_{y\in Y}\mathsf C(y)
                         \prod_{u\in U}\mathsf B(u)|0\rangle.
\end{equation}
It is not defined using a Hermitian adjoint; that identification will
be proved separately.  Write $C(Y,U)_{jk}=1/[U_k-y_j]$, and put
\begin{align}
 m_u^t(y)&=-\frac1{2\pi i}\bigl[(1+e^{-2it}q_Y(u))\coth(u-y)
                 \notag\\
 &\hspace{24mm}-\coth(u-y-i\lambda)
                 -e^{-2it}q_Y(u)\coth(u-y+i\lambda)\bigr],
 \label{eq:finite-source-column}\\
 M^t(Y,U)_{jk}&=m_{U_k}^t(y_j).
 \label{eq:reduced-scalar-matrix}
\end{align}
All root coincidences in these matrices mean analytic limits along
the full formula.  In particular
\begin{equation}\label{eq:source-Gaudin-limit}
 M^t(Y,Y)=G_Y.
\end{equation}

The derivatives $\partial_{y_j}\tau_t$ below are taken with respect
to the independent root coordinates before imposing the on-shell
equations; they are not derivatives along a Bethe-root branch.

\begin{lemma}[Exact scalar products and deleted columns]
\label{lem:scalar-product}
In the normalization \eqref{eq:standard-monodromy},
\begin{equation}\label{eq:Slavnov-exact}
 S(Y,U)=e^{-int}\left(\prod_{y\in Y}l(y)\right)
 \frac{\det_{j,k}\partial_{y_j}\tau_t(U_k;Y)}
      {\det C(Y,U)}.
\end{equation}
Consequently, if $\det G_Y\ne0$, then $S(Y,Y)\ne0$ and
\begin{equation}\label{eq:Gaudin-norm-exact}
 S(Y,Y)=(2\pi i s_\lambda)^n
       \left(\prod_{y\in Y}a(y)l(y)\right)
       \left(\prod_{j\ne k}f(y_j-y_k)\right)\det G_Y.
\end{equation}
Suppose $U$ is obtained by replacing the roots indexed by
$H=(p_1,\ldots,p_k)$ by $V=(v'_1,\ldots,v'_k)$, in those same slots.
Set
\[
 A_Y(p)=a(p)\prod_{y\in Y\setminus\{p\}}f(y-p)\quad(p\in Y).
\]
Then the normalized scalar product is exactly
\begin{equation}\label{eq:deleted-scalar-product}
 \frac{S(Y,U)}{S(Y,Y)}
 =(-s_\lambda)^{-k}
   \frac{\prod_{v'\in V}a(v')F_Y(v')}{\prod_{p\in H}A_Y(p)}
   \frac{\det\bigl((G_Y^{-1}m_{v'}^t)_p\bigr)_{H,V}}
        {\det\bigl(1/[v'-p]\bigr)_{H,V}}.
\end{equation}
If the positive weights \eqref{eq:Gaudin-weights} are defined, the
last numerator equals
\begin{equation}\label{eq:W-determinant}
 \det\bigl((G_Y^{-1}m_{v'}^t)_p\bigr)_{H,V}
 =\left(\prod_{p\in H}\ell_p\right)
       \det\bigl(W_{v'}(p)\bigr)_{H,V},
 \qquad W_{v'}=(I+K\ell)^{-1}m_{v'}^t.
\end{equation}
Empty products and determinants have value one.
\end{lemma}

\begin{proof}
We apply the trigonometric scalar-product theorem to the fundamental
six-vertex monodromy \eqref{eq:standard-monodromy}, with highest vector
$|0\rangle$, vacuum functions $a,l$, and diagonal twist
$\kappa=e^{-2it}$. To match the conventions of
\cite{KMST2005}, take its anisotropy to be $i\lambda$ and its
inhomogeneities to be $b_j+i\lambda/2$, with a relabeling of sites
if necessary. Its monodromy is $s_\lambda^L M(x)$, so the scalar
product and the vacuum-times-derivative determinant both acquire the
same factor $s_\lambda^{2Ln}$. After cancelling this factor, the
prefactor is $\prod_{y\in Y}l(y)$ for the eigenvalue
$\mathsf A+\kappa\mathsf L$; this is
\cite[Proposition~3.1, equations~(3.8)--(3.9)]{KMST2005}.
Replacing that eigenvalue by $\tau_t=e^{it}\tau_\kappa$ multiplies its
$n$ derivative columns by $e^{it}$ and therefore gives the compensating
factor $e^{-int}$ in \eqref{eq:Slavnov-exact}.  The untwisted
derivative-over-Cauchy version, with its equal-spin-normalized
monodromy, is \cite[Theorem~3.1, equations~(3.4)--(3.5)]{KMT1999}.
These are finite scalar-product identities; they make no assertion
about completeness or about a thermodynamic limit.

Differentiation of \eqref{eq:bethe-eigenvalue} gives
\[
 \partial_{y_j}\tau_t(u;Y)
     =e^{it}a(u)F_Y(u)(-2\pi i)m_u^t(y_j).
\]
The factor $e^{int}$ cancels that in
\eqref{eq:Slavnov-exact}.  When $u\to y_k$, the product
$(1+e^{-2it}q_Y(u))\coth(u-y_k)$ tends to
$-2\pi i\xi_Y'(y_k)$, because
$q_Y'(y_k)=-2\pi i e^{2it}\xi_Y'(y_k)$.
The two remaining coth terms give $K(y_k-y_j)$ after division by
$-2\pi i$.  This proves \eqref{eq:source-Gaudin-limit}, including its
diagonal.

At a root $p$ the residue of the extracted vacuum factor is
\begin{equation}\label{eq:vacuum-residue}
 \lim_{u\to p}[u-p]a(u)F_Y(u)=-s_\lambda A_Y(p).
\end{equation}
Cancel these poles with the corresponding Cauchy poles.  If every
column tends to its root, the remaining determinant is $\det G_Y$,
giving \eqref{eq:Gaudin-norm-exact}.  For partial coincidence first
cancel the unchanged columns.  The determinant of the matrix obtained
from $G_Y$ by replacing columns $H$ by $m_{v'}^t$, divided by
$\det G_Y$, is
$\det((G_Y^{-1}m_{v'}^t)_p)_{H,V}$: multiply the matrix on the left by
$G_Y^{-1}$ and expand along its unchanged unit columns.  Expanding the
Cauchy determinant along its unchanged poles leaves exactly
$\det(1/[v'-p])_{H,V}$ with the same slot order.  Each deleted root
loses the factor in \eqref{eq:vacuum-residue}.  This proves
\eqref{eq:deleted-scalar-product}, including its sign and constant.
Finally $G_Y^{-1}=\diag(\ell)(I+K\ell)^{-1}$ proves
\eqref{eq:W-determinant}.

All factors in the norm outside $\det G_Y$ are nonzero: real $y,b$
give $[y-b\pm i\lambda/2]\ne0$, and distinct real roots give
$[y-y']\ne0$ and $[y-y'\pm i\lambda]\ne0$.
The scalar-product theorem is first used at generic variables.
At any regular on-shell specialization, the nonsingular counting
Jacobian permits a local analytic perturbation of the roots along
with the parameters.  Meromorphic continuation and the pole
cancellations above then give the asserted formulas there as well.
In particular, no exclusion at a root of unity is necessary.
\end{proof}

\subsection{Hilbert duality and the finite-word formula}

\begin{lemma}[The Hilbert dual]\label{lem:hilbert-dual}
For even $L$ and real inhomogeneities,
\begin{equation}\label{eq:B-Hilbert-adjoint}
 \mathsf B(x)^*=-\mathsf C(\bar x).
\end{equation}
Thus for a distinct real on-shell list with $\det G_Y\ne0$, the
physical unit vector
\begin{equation}\label{eq:physical-Bethe-vector}
 \Omega_Y=\frac{\mathcal P_L|Y\rangle}{\|\mathcal P_L|Y\rangle\|}
\end{equation}
is well defined, and
\[
 \|\mathcal P_L|Y\rangle\|^2=(-1)^nS(Y,Y)>0.
\]
If $Y,Z$ are two such real on-shell lists of the same size, possibly
for different twists, then
\begin{equation}\label{eq:two-orientation-overlap}
 |\langle\Omega_Y,\Omega_Z\rangle|^2
       =\frac{S(Y,Z)S(Z,Y)}{S(Y,Y)S(Z,Z)}.
\end{equation}
These identities also hold for real staggered $b_j$, without a Perron
interpretation of the vectors.
\end{lemma}

\begin{proof}
Polarize the adjoint identity in Lemma~\ref{lem:row-algebra} over the
four auxiliary marks.  It gives
$\mathsf B_{\rm ph}(x)^*=\mathsf C_{\rm ph}(\bar x)$.
Using the gauge identity and commuting $\mathcal P_L$ with
$\mathcal Q_L$ gives
$\mathsf B(x)^*=\mathcal Q_L\mathsf C(\bar x)\mathcal Q_L$.
Since $\mathsf C$ changes the number of minus spins by one, it
anticommutes with $\mathcal Q_L$, proving
\eqref{eq:B-Hilbert-adjoint}.  Creation entries commute, so for real
roots
\[
 (\mathcal P_L|Y\rangle)^*
       =(-1)^n\langle0|\prod_{y\in Y}\mathsf C(y)\mathcal P_L.
\]
Lemma~\ref{lem:scalar-product} makes this norm nonzero, and positivity
is then Hilbert-space positivity.  Applying the same identity to both
cross scalar products proves \eqref{eq:two-orientation-overlap}.
Absorbing the twist into the monodromy merely replaces
$\mathsf B$ by $e^{it}\mathsf B$ and $\mathsf C$ by
$e^{-it}\mathsf C$, preserving the adjoint relation and the normalized
ratio.
\end{proof}

We can now combine the entry actions with the scalar-product formula.
The result is a sum over the original roots removed by a row word.
Its determinant size depends only on the word length, not on the chain
length; this is the form needed for the thermodynamic limit.  In the remainder of this subsection $t=0$, $Y$ is a
regular real on-shell list, and the spectral arguments $v'_1,\ldots,v'_m$
are distinct and generic.  Define normalized physical rows by
\begin{equation}\label{eq:normalized-finite-row}
 A_L^D(x)=\frac{T^D(x)}{(-1)^n\tau_0(x;Y)},\qquad
 P_L=A_L^I,\qquad J_L=A_L^Z,
\end{equation}
on the open set where the denominator is nonzero.  Then
$P_L(x)\Omega_Y=\Omega_Y$.  Holomorphic continuation and norm bounds
for these rows will be proved using the estimates of
Section~\ref{sec:thermodynamic}.

For each path in Lemma~\ref{lem:finite-actions}, let $T$ be the set
remaining just after its current removal and define
\begin{align}
 U_Y(T,p)&=
  \frac{\prod_{u\in T\setminus Y}f(u-p)}
       {\prod_{y\in Y\setminus(T\cup\{p\})}f(y-p)},&
 L_Y(T,p)&=
  \frac{\prod_{u\in T\setminus Y}f(p-u)}
       {\prod_{y\in Y\setminus(T\cup\{p\})}f(p-y)},
 \label{eq:path-quotients}\\
 \chi_Y(p)&=
  \begin{cases}1,&p\in Y,\\q_Y(p),&p\notin Y.\end{cases}
 \label{eq:path-lower-factor}
\end{align}
Associate with an action at $v'$ the correction factor
\begin{equation}\label{eq:path-correction}
 \gamma=
 \begin{cases}
  1,&\mathsf B\text{ (no removal)},\\
  \beta(v'-p)U_Y(T,p),&\mathsf A\text{ (remove }p),\\
  \chi_Y(p)\beta(p-v')L_Y(T,p),&\mathsf L\text{ (remove }p),\\
  \chi_Y(q)\beta(v'-p)\beta(q-v')f(q-p)
        U_Y(T,p)L_Y(T,q),&\mathsf C\text{ (remove }p,q).
 \end{cases}
\end{equation}
These are explicit rational functions of hyperbolic sines.  The only
dependence on the original root set outside the factors displayed in
\eqref{eq:path-quotients} is through $q_Y$ at added arguments.

\begin{corollary}[Exact finite-word expansion]\label{cor:finite-word-expansion}
Expand each auxiliary mark as
\[
 T^{D_j}=\sum_{\alpha,\beta}(D_j)_{\beta\alpha}
 M_{{\rm ph},\alpha\beta}.
\]
For each resulting entry word, let $n_j$
be the number of minus spins after its $j$th action, with $n_0=n$.
Entry words with $n_m\ne n$ have zero expectation in $\Omega_Y$.
For a path of a word with $n_m=n$, let $H\subset Y$ be the roots
absent at the end and let $V$ be its surviving added arguments, ordered
in the same replacement slots.  Then $|H|=|V|=:k\le m$, and its
contribution to the normalized physical expectation is
\begin{equation}\label{eq:exact-path-determinant}
 \left(\prod_{j=1}^m(D_j)_{\beta_j\alpha_j}\right)
 (-1)^{\sum_{j=1}^m(n_{j-1}-n)}
 (-s_\lambda)^{-k}
 \frac{\prod_{j=1}^m\gamma_j}
      {\prod_{j=1}^m(1+q_Y(v'_j))}
 \frac{\det(W_{v'}(p))_{H,V}}
      {\det(1/[v'-p])_{H,V}}
 \prod_{p\in H}\ell_p.
\end{equation}
The expectation is the sum of these contributions.  In particular,
the determinant size is bounded by the word length.  For diagonal
entry words the parity sign in \eqref{eq:exact-path-determinant} is
one.  Among distinct removed original roots the correction product
contains only reciprocal $f$ factors, with one for each pair.  Each
removed original root also contributes one factor $\beta$ joining it
to the row argument at the step that removes it.
\end{corollary}

\begin{proof}
For this proof extend the vacuum base to added arguments by
\[
 \mathcal A_Y(p)=
 \begin{cases}
  A_Y(p),&p\in Y,\\
  a(p)F_Y(p),&p\notin Y.
 \end{cases}
\]
The coefficient of an upper removal is
$\mathcal A_Y(p)\beta(v'-p)U_Y(T,p)$.
The corresponding lower coefficient is
$\mathcal A_Y(p)\chi_Y(p)\beta(p-v')L_Y(T,p)$:
on a root the equality of the two vacuum bases is precisely
\eqref{eq:bethe-equations} at $t=0$, including its self factor $-1$;
off a root their ratio is $q_Y(p)$.
For two removals the extra factor is $f(q-p)$, exactly as in
\eqref{eq:C-action}.  Thus the coefficient of a path before taking its
scalar product is the product of its corrections $\gamma_j$ times one
$\mathcal A_Y(p)$ for every removed element.

Every original root is either retained or removed once.  Similarly,
each new spectral argument is either retained in $V$ or removed once,
possibly in the very action that introduced it.  Consequently the
product of vacuum bases along the path is
\[
 \left(\prod_{p\in H}A_Y(p)\right)
 \left(\prod_{v'\in\{v'_1,\ldots,v'_m\}\setminus V}
                       a(v')F_Y(v')\right).
\]
Multiply by \eqref{eq:deleted-scalar-product}.  The deleted-root bases
cancel, and the surviving arguments supply the missing vacuum bases.
Dividing by $\prod_j\tau_0(v'_j;Y)$ leaves exactly
$\prod_j(1+q_Y(v'_j))^{-1}$ and $(-s_\lambda)^{-k}$.
Equation~\eqref{eq:W-determinant} supplies the determinant and weights.

The Hilbert dual contributes the same $(-1)^n$ to the scalar product
and norm, so it cancels by Lemma~\ref{lem:hilbert-dual}.  Each physical
entry, after division by $(-1)^n$, contributes
$(-1)^{n_{j-1}-n}$ from $\mathcal Q_L$ on its input sector; adjacent
$\mathcal P_L$ factors cancel.  This proves the explicit parity sign.
Different final spin sectors are orthogonal, proving the zero
assertion.

Finally consider a pair of deleted original roots.  If removed at
different steps, the earlier root occurs in the denominator of
\eqref{eq:path-quotients} at the later removal.  If removed together,
the two denominators contribute $f(q-p)^{-2}$ and the extra factor
$f(q-p)$ leaves one reciprocal.  There are no other factors involving
two original roots.  The factor $\beta$ at each original removal is
explicit in \eqref{eq:path-correction}.  This proves the last claims.
\end{proof}

For real $u$, $|f(u)|^2=1+\sin^2\lambda/\sinh^2u\ge1$.
Thus the reciprocal factors between deleted roots cause no collision
growth, and for a bounded generic added argument the corresponding
$\beta$ factor decays like $e^{-|p|}$ as its real root $p$ escapes.
These bounds control the pair factors and the individual removal factors.
Simultaneous escape of several roots still requires a determinant bound,
which will be proved in Section~\ref{sec:thermodynamic}.

\subsection{Two-twist overlaps}

The normalization argument in Section~\ref{sec:normalization} will
also use a different consequence of the same scalar-product formula:
the overlap of the unit vectors associated with two twists.
Unlike a fixed row word, this overlap retains determinants of the full
root-set size. We record its exact pair product here, before taking
any thermodynamic limit.

To factor \eqref{eq:two-orientation-overlap}, denote the twists of $Y,Z$ by
$t_Y,t_Z$, respectively.  Put
\[
 \Delta(Y)=\prod_{j<k}[y_k-y_j],\qquad
 \mathcal H(Y,U)=\prod_{j,k}[y_j-U_k+i\lambda],\qquad
 c_n=(2\pi i)^n(-1)^{n(n-1)/2}.
\]
The trigonometric Cauchy identity rewrites
\eqref{eq:Slavnov-exact} as
\begin{equation}\label{eq:Slavnov-separated}
 S(Y,U)=c_n\left(\prod_{y\in Y}l(y)\right)
             \left(\prod_{u\in U}a(u)\right)
       \frac{\mathcal H(Y,U)}{\Delta(Y)\Delta(U)}
             \det M^{t_Y}(Y,U).
\end{equation}
Consequently all one-set factors cancel separately, leaving
\begin{equation}\label{eq:overlap-determinant-finite}
 |\langle\Omega_Y,\Omega_Z\rangle|^2
 =\frac{\mathcal H(Y,Z)\mathcal H(Z,Y)}
        {\mathcal H(Y,Y)\mathcal H(Z,Z)}
   \frac{\det M^{t_Y}(Y,Z)\det M^{t_Z}(Z,Y)}
        {\det G_Y\det G_Z}.
\end{equation}
This is a finite identity, including the full shifted pair product.
In particular, along a finite regular twist branch $Y(t)$, take
$Y=Y(0)$ and $Z=Y(t)$. The overlap is nonzero for $t$ near zero,
since its value at zero is one.  Its local logarithm and derivatives
there are therefore well defined.  Neither changing the phases of the
unit vectors nor changing the logarithm branch changes the quadratic
coefficient.  No neighborhood uniform in $L$ is asserted here.


We now return to the fixed-word formula
\eqref{eq:exact-path-determinant}. Its scalar coefficients are explicit,
but passing to the plane requires uniform control of the root measures,
the Gaudin inverses, and determinants with several roots in the sparse
ends. Section~\ref{sec:thermodynamic} supplies those estimates.
